% ch26_ends_and_coends.tex — Volume III Chapter 2

\chapter{Ends and Coends}

\begin{overviewbox}
Volume II made repeated unsupported promises about \emph{ends} and
\emph{coends}: profunctor composition (Chapter~14) used the formula
$(Q \circ P)(c,a) = \int^b Q(c,b) \times P(b,a)$ without explaining the
integral sign; the Kan extension colimit formula
$(\operatorname{Lan}_p F)(e) = \int^c \mathcal{E}(p(c), e) \cdot F(c)$ likewise
sidestepped what $\int^c$ means; the enriched Yoneda lemma (Chapter~20)
even \emph{stated} its conclusion as an end $\int_c \mathcal{V}(G(c), H(c))$.

This chapter delivers on those promises. Ends and coends are the universal
constructions for ``mixed-variance'' functors $T : \mathcal{C}^{\mathrm{op}}
\times \mathcal{C} \to \mathcal{D}$ — limits or colimits that, on the diagonal,
``contract'' the two copies of $\mathcal{C}$ into one. They are the
foundational vocabulary of formal category theory: weighted limits, enriched
Kan extensions, Day convolution, and the formal Yoneda lemma are all
expressed in coend calculus. After this chapter, the rest of Volume III
flows.
\end{overviewbox}


% ═══════════════════════════════════════════════════════════════════════════
\section{Dinatural Transformations}
\label{sec:dinatural}
% ═══════════════════════════════════════════════════════════════════════════

\subsection{Formalizing}

Let $T : \mathcal{C}^{\mathrm{op}} \times \mathcal{C} \to \mathcal{D}$ be a
functor — a \emph{mixed-variance} or \emph{bifunctor} in the contravariant-
then-covariant sense. The objects of interest are the diagonal values
$T(c, c)$ and how they fit together across $\mathcal{C}$.

\begin{definition}[Dinatural transformation]
\label{def:dinatural}
A \term{dinatural transformation} $\alpha : T \xRightarrow{..} S$ between two
such functors $T, S : \mathcal{C}^{\mathrm{op}} \times \mathcal{C} \to
\mathcal{D}$ is a family of morphisms $\alpha_c : T(c,c) \to S(c,c)$,
one for each $c \in \mathcal{C}$, such that for every $f : c \to c'$ the
hexagon
\begin{equation*}
\begin{tikzcd}[column sep=tiny, row sep=small]
  & T(c, c) \arrow[r, "\alpha_c"] & S(c, c) \arrow[dr, "{S(\mathrm{id}, f)}"] & \\
  T(c', c) \arrow[ur, "{T(f, \mathrm{id})}"] \arrow[dr, "{T(\mathrm{id}, f)}" '] & & & S(c, c') \\
  & T(c', c') \arrow[r, "\alpha_{c'}" '] & S(c', c') \arrow[ur, "{S(f, \mathrm{id})}" '] &
\end{tikzcd}
\end{equation*}
commutes. We call this the \term{hexagonal} or \term{dinaturality} condition.
\end{definition}

\begin{howtosay}
$\alpha : T \xRightarrow{..} S$ \sayit{``dinatural alpha from T to S''}, with
the dotted arrow signaling that the naturality holds in the hexagonal
(dinatural) sense rather than the square (natural) sense.
\end{howtosay}

\begin{remark}[When the hexagon collapses to a square]
If $T$ and $S$ are constant in their first argument — i.e., they factor
through the projection $\mathcal{C}^{\mathrm{op}} \times \mathcal{C} \to
\mathcal{C}$ — the hexagon collapses to an ordinary naturality square. Thus
ordinary natural transformations are the special case of dinatural
transformations between covariant functors. The dinatural notion is strictly
more general.
\end{remark}

\subsection{Reorganizing: Wedges as Dinatural Cones}

\begin{definition}[Wedge]
\label{def:wedge}
A \term{wedge} from an object $d \in \mathcal{D}$ to a functor
$T : \mathcal{C}^{\mathrm{op}} \times \mathcal{C} \to \mathcal{D}$ is a
dinatural transformation $\omega : \Delta d \xRightarrow{..} T$ from the
constant functor at $d$. Concretely: a family $\omega_c : d \to T(c, c)$
such that for every $f : c \to c'$ the square
\begin{equation*}
\begin{tikzcd}
  d \arrow[r, "\omega_{c'}"] \arrow[d, "\omega_c" '] & T(c', c') \arrow[d, "{T(f, \mathrm{id})}"] \\
  T(c, c) \arrow[r, "{T(\mathrm{id}, f)}" '] & T(c, c')
\end{tikzcd}
\end{equation*}
commutes. (The hexagon for $\Delta d$ collapses because $\Delta d$ has no
non-trivial action.) Dually, a \term{cowedge} from $T$ to $d$ is a dinatural
$\nu : T \xRightarrow{..} \Delta d$, i.e., a family $\nu_c : T(c, c) \to d$
satisfying the analogous square.
\end{definition}

A wedge replaces ``a cone from $d$'' (in the limit picture) and a cowedge
replaces ``a cone over $d$'' (in the colimit picture), now for the
mixed-variance functor $T$.


% ═══════════════════════════════════════════════════════════════════════════
\section{The End}
\label{sec:end}
% ═══════════════════════════════════════════════════════════════════════════

\subsection{Formalizing}

\begin{definition}[End]
\label{def:end}
Let $T : \mathcal{C}^{\mathrm{op}} \times \mathcal{C} \to \mathcal{D}$. An
\term{end} of $T$ is an object $\int_c T(c, c) \in \mathcal{D}$ together with
a universal wedge $\pi : \Delta \int_c T(c,c) \xRightarrow{..} T$: for every
wedge $\omega : \Delta d \xRightarrow{..} T$ there exists a unique morphism
$\hat\omega : d \to \int_c T(c,c)$ with $\pi_c \circ \hat\omega = \omega_c$
for all $c$.
\end{definition}

\begin{howtosay}
$\int_c T(c, c)$ \sayit{``the end over c of T(c, c)''} or just \sayit{``end of
T''}. The subscript $c$ on the integral sign marks the bound variable; it does
not denote a particular value of $c$.
\end{howtosay}

\subsection{Reorganizing: Ends as Equalizers}

\begin{theorem}[End as an equalizer]
\label{thm:end-as-equalizer}
The end $\int_c T(c, c)$, when $\mathcal{D}$ has products and equalizers and
$\mathcal{C}$ is small, exists and is computed as the equalizer
\begin{equation*}
  \int_c T(c, c)
  \;\xrightarrow{e}\;
  \prod_c T(c, c)
  \;\rightrightarrows\;
  \prod_{f : c \to c'} T(c, c')
\end{equation*}
where the two parallel arrows are induced respectively by $T(\mathrm{id}, f)$
and $T(f, \mathrm{id})$, indexed over all morphisms $f$ of $\mathcal{C}$.
\end{theorem}

\begin{prooflog}
\proofstep{An element of $\prod_c T(c,c)$ is a family $(\omega_c)_c$ with
$\omega_c \in T(c,c)$.}{Universal property of the product.}
\proofstep{The two parallel arrows send $(\omega_c)$ to the families
$(T(\mathrm{id}, f) \omega_c)_{f}$ and $(T(f, \mathrm{id}) \omega_{c'})_{f}$
in $\prod_{f} T(c, c')$.}{By definition of the induced maps into the
product over morphisms.}
\proofstep{The equalizer condition $T(\mathrm{id}, f) \omega_c = T(f,
\mathrm{id}) \omega_{c'}$ for every $f : c \to c'$ is precisely the wedge
condition of Definition~\ref{def:wedge}.}{Direct comparison of the equations.}
\proofstep{So the equalizer is the universal object equipped with a wedge
$\pi_c : E \to T(c,c)$ for each $c$, satisfying exactly the wedge
condition.}{Equalizer universal property.}
\proofstep{Therefore the equalizer is the end.}{By Definition~\ref{def:end}.}
\end{prooflog}

\begin{corollary}
Ends exist whenever $\mathcal{D}$ is complete (has small limits) and
$\mathcal{C}$ is small. They are limits — limits of a specific diagram of
shape determined by $T$.
\end{corollary}

\subsection{Formally: Examples}

\begin{example_thm}[Natural transformations as an end]
\label{ex:nat-as-end}
For functors $F, G : \mathcal{C} \to \mathcal{D}$, define
$T(c, c') = \mathcal{D}(F c, G c')$. Then
\begin{equation*}
  \int_c \mathcal{D}(Fc, Gc) \;\cong\; \mathrm{Nat}(F, G).
\end{equation*}
The end's wedge $\pi_c : \mathrm{Nat}(F, G) \to \mathcal{D}(Fc, Gc)$ is
evaluation at $c$; the wedge condition is exactly the naturality square.
\end{example_thm}

\begin{example_thm}[Enriched natural transformations]
For $\mathcal{V}$-functors $G, H : \mathcal{C} \to \mathcal{V}$ between
$\mathcal{V}$-categories, the object of $\mathcal{V}$-natural transformations
is the end
\begin{equation*}
  [\mathcal{C}, \mathcal{V}]_\mathcal{V}(G, H) \;=\; \int_c \mathcal{V}(G c, H c)
\end{equation*}
in $\mathcal{V}$ (cf.\ Chapter~20). This is the enriched generalization of
Example~\ref{ex:nat-as-end}.
\end{example_thm}

\begin{example_thm}[Center of a category]
For a category $\mathcal{C}$ viewed as a one-object 2-category, the end
$\int_c \mathcal{C}(c, c)$ over $\mathcal{C}$ itself (with $T(c, c') =
\mathcal{C}(c, c')$) is the monoid of \term{natural endomorphisms} of the
identity functor — the \emph{center} of $\mathcal{C}$.
\end{example_thm}


% ═══════════════════════════════════════════════════════════════════════════
\section{The Coend}
\label{sec:coend}
% ═══════════════════════════════════════════════════════════════════════════

\subsection{Formalizing}

\begin{definition}[Coend]
\label{def:coend}
The \term{coend} $\int^c T(c, c) \in \mathcal{D}$ of a functor $T :
\mathcal{C}^{\mathrm{op}} \times \mathcal{C} \to \mathcal{D}$ is an object
equipped with a universal cowedge $\iota : T \xRightarrow{..} \Delta \int^c
T(c, c)$: for every cowedge $\nu : T \xRightarrow{..} \Delta d$ there is a
unique morphism $\hat\nu : \int^c T(c,c) \to d$ with $\hat\nu \circ \iota_c
= \nu_c$ for all $c$.
\end{definition}

\begin{theorem}[Coend as a coequalizer]
\label{thm:coend-as-coequalizer}
When $\mathcal{D}$ has coproducts and coequalizers and $\mathcal{C}$ is small,
the coend exists and is the coequalizer
\begin{equation*}
  \coprod_{f : c \to c'} T(c', c) \;\rightrightarrows\; \coprod_c T(c, c)
  \;\xrightarrow{\;\iota\;}\; \int^c T(c, c).
\end{equation*}
The two parallel arrows are induced by $T(f, \mathrm{id}) : T(c', c) \to
T(c, c)$ and $T(\mathrm{id}, f) : T(c', c) \to T(c', c')$.
\end{theorem}

\begin{proof}
Dual to Theorem~\ref{thm:end-as-equalizer}, with products replaced by
coproducts, equalizers by coequalizers, and the universal property reversed.
A cowedge from $T$ to $d$ is precisely a family $(\nu_c)$ satisfying
$\nu_c \circ T(f, \mathrm{id}) = \nu_{c'} \circ T(\mathrm{id}, f)$ for each
$f$, which is the coequalizer condition.
\end{proof}

\subsection{Reorganizing: The Tensor-Product Picture}

The coend $\int^c T(c, c)$ admits a useful informal reading as a \emph{tensor
product}. For $\mathcal{D} = \mathbf{Set}$ and $T(c, c') = P(c) \times Q(c')$
with $P$ contravariant and $Q$ covariant,
\begin{equation*}
  \int^c P(c) \times Q(c)
  \;\cong\;
  \frac{\coprod_c P(c) \times Q(c)}{\sim}
\end{equation*}
where $\sim$ identifies $(P(f)(p), q) \sim (p, Q(f)(q))$ for $f : c \to c'$.
This is the standard quotient defining a tensor product. In general,
``$\int^c$'' is the tensor product over $\mathcal{C}$ in the same sense that
tensor products of modules are coequalizers over a ring.

\subsection{Formally: Examples}

\begin{example_thm}[Profunctor composition]
\label{ex:profunctor-composition}
For profunctors $P : \mathcal{A} \to \mathcal{B}$ (i.e., $P :
\mathcal{B}^{\mathrm{op}} \times \mathcal{A} \to \mathbf{Set}$) and $Q :
\mathcal{B} \to \mathcal{C}$, their composition is
\begin{equation*}
  (Q \circ P)(c, a) \;=\; \int^{b \in \mathcal{B}} Q(c, b) \times P(b, a).
\end{equation*}
This is the formula promised in Chapter~14. The coend identifies pairs
$(P(f)(p), q) \sim (p, Q(f)(q))$ along the morphism $f : b \to b'$, the
correct ``mixing'' of $P$'s output with $Q$'s input.
\end{example_thm}

\begin{example_thm}[Left Kan extension]
\label{ex:lan-coend}
For $F : \mathcal{C} \to \mathcal{E}$ and $K : \mathcal{C} \to \mathcal{D}$,
\begin{equation*}
  (\operatorname{Lan}_K F)(d) \;\cong\; \int^c \mathcal{D}(Kc, d) \cdot F(c)
\end{equation*}
where $S \cdot X$ for $S \in \mathbf{Set}$ and $X \in \mathcal{E}$ denotes
the copower $\coprod_{S} X$. This is the formula from Chapter~18, now with
the integral sign defined.
\end{example_thm}

\begin{example_thm}[Geometric realization]
The geometric realization of a simplicial set $X : \Delta^{\mathrm{op}} \to
\mathbf{Set}$ is
\begin{equation*}
  |X| \;=\; \int^{[n] \in \Delta} X_n \cdot \Delta^n
\end{equation*}
where $\Delta^n$ is the topological $n$-simplex and $\cdot$ is the topological
copower $S \cdot Y = \coprod_S Y$. The coend identifies face and degeneracy
maps with their topological counterparts.
\end{example_thm}


% ═══════════════════════════════════════════════════════════════════════════
\section{The Coend Calculus}
\label{sec:coend-calculus}
% ═══════════════════════════════════════════════════════════════════════════

The reason ends and coends earn the integral notation is that they obey
calculus-like rules: a Fubini theorem for swapping the order of integration,
a substitution rule, and a Yoneda reduction. These rules turn coend
manipulations into mechanical computations.

\subsection{Fubini}

\begin{theorem}[Fubini for ends]
\label{thm:fubini-end}
For a functor $T : (\mathcal{C} \times \mathcal{D})^{\mathrm{op}} \times
(\mathcal{C} \times \mathcal{D}) \to \mathcal{E}$,
\begin{equation*}
  \int_{(c, d)} T((c, d), (c, d))
  \;\cong\;
  \int_c \int_d T((c, d), (c, d))
  \;\cong\;
  \int_d \int_c T((c, d), (c, d))
\end{equation*}
whenever the relevant limits exist.
\end{theorem}

\begin{theorem}[Fubini for coends]
\label{thm:fubini-coend}
Dually,
\begin{equation*}
  \int^{(c, d)} T((c, d), (c, d))
  \;\cong\;
  \int^c \int^d T((c, d), (c, d))
  \;\cong\;
  \int^d \int^c T((c, d), (c, d)).
\end{equation*}
\end{theorem}

\begin{proof}[Sketch]
Both statements follow from the fact that ends are limits and coends are
colimits, and limits/colimits over a product category factor through limits/
colimits over each factor (Fubini for ordinary limits). The end of a
mixed-variance functor uses the same Fubini at the level of equalizer
diagrams.
\end{proof}

\subsection{Yoneda Reduction}

\begin{theorem}[Yoneda reduction — coend form]
\label{thm:yoneda-reduction-coend}
For a (covariant) functor $F : \mathcal{C} \to \mathbf{Set}$,
\begin{equation*}
  \int^c \mathcal{C}(c, c') \times F(c) \;\cong\; F(c').
\end{equation*}
Dually for $G : \mathcal{C}^{\mathrm{op}} \to \mathbf{Set}$,
\begin{equation*}
  \int^c G(c) \times \mathcal{C}(c', c) \;\cong\; G(c').
\end{equation*}
\end{theorem}

\begin{howtosay}
This is called the \term{ninja Yoneda lemma}: the variable $c$ is
``integrated out'' against the representable $\mathcal{C}(c, c')$, leaving
only the value at $c'$. The name comes from how silently and decisively the
coend disappears.
\end{howtosay}

\begin{prooflog}
\proofstep{Compute the coend explicitly as the coequalizer
$\coprod_{f : c \to d} \mathcal{C}(d, c') \times F(c) \rightrightarrows
\coprod_c \mathcal{C}(c, c') \times F(c)$.}{By Theorem~\ref{thm:coend-as-coequalizer}.}
\proofstep{The two parallel arrows identify $(\mathcal{C}(d, c') \ni g, F(c)
\ni x) \sim (g \circ f, x)$ vs.\ $(g, F(f)(x))$ along $f : c \to d$.}{Direct
substitution.}
\proofstep{Define $\phi : \coprod_c \mathcal{C}(c, c') \times F(c) \to
F(c')$ by $(g, x) \mapsto F(g)(x)$.}{The candidate map to $F(c')$.}
\proofstep{$\phi$ respects the equivalence relation: $\phi(g \circ f, x)
= F(g \circ f)(x) = F(g)(F(f)(x)) = \phi(g, F(f)(x))$.}{By functoriality of
$F$.}
\proofstep{Therefore $\phi$ factors through the coend to give $\hat\phi :
\int^c \mathcal{C}(c, c') \times F(c) \to F(c')$.}{By the universal property
of the coequalizer.}
\proofstep{Construct the inverse $\psi : F(c') \to \int^c \mathcal{C}(c, c')
\times F(c)$ by $x \mapsto [\mathrm{id}_{c'}, x]$ (the equivalence class).}{
The standard ``identity-and-itself'' element.}
\proofstep{$\hat\phi \circ \psi = \mathrm{id}_{F(c')}$: $\hat\phi([\mathrm{id},
x]) = F(\mathrm{id})(x) = x$.}{Identity action.}
\proofstep{$\psi \circ \hat\phi = \mathrm{id}$ on classes: for $(g, x)$ with
$g : c \to c'$, $\psi(\phi(g, x)) = [\mathrm{id}_{c'}, F(g)(x)] \sim
[g, x]$ via the relation at $f = g$.}{The relation collapses every class to
a singleton with identity.}
\proofstep{Both composites are the identity; isomorphism established.}{
Yoneda reduction.}
\end{prooflog}

\begin{remark}[Why this is Yoneda]
Theorem~\ref{thm:yoneda-reduction-coend} is the Yoneda lemma in disguise.
The ordinary Yoneda lemma says elements of $F(c')$ correspond to natural
transformations $\mathcal{C}(c', -) \Rightarrow F$, which are an end:
$F(c') \cong \int_c \mathbf{Set}(\mathcal{C}(c', c), F(c))$. Taking the
hom-set adjunction backwards (using that the coend is the colimit and the
end is the limit) gives the coend formulation. The two forms are duals:
ends pair with hom-functors covariantly; coends pair with hom-functors
contravariantly.
\end{remark}

\subsection{The Co-Yoneda Lemma}

\begin{theorem}[Co-Yoneda / density]
\label{thm:co-yoneda}
For any functor $F : \mathcal{C} \to \mathbf{Set}$,
\begin{equation*}
  F \;\cong\; \int^c \mathcal{C}(c, -) \cdot F(c)
\end{equation*}
as functors $\mathcal{C} \to \mathbf{Set}$. Every functor is a colimit of
representables.
\end{theorem}

\begin{proof}
At each $c' \in \mathcal{C}$, applying the formula gives $\int^c
\mathcal{C}(c, c') \times F(c) \cong F(c')$ by
Theorem~\ref{thm:yoneda-reduction-coend}. Naturality in $c'$ is automatic
because the coend is functorial in its non-integrated arguments.
\end{proof}

This is the precise statement that ``the presheaf category is the free
cocompletion of $\mathcal{C}$'' (Chapter~18, Chapter~24): every presheaf is
canonically a colimit of representables, witnessed by the coend formula.


% ═══════════════════════════════════════════════════════════════════════════
\section{Ends and Coends in Enriched Settings}
\label{sec:enriched-ends-coends}
% ═══════════════════════════════════════════════════════════════════════════

\subsection{Formalizing}

The development so far has assumed $\mathcal{D}$ is an ordinary category. In
the enriched setting, $\mathcal{D} = \mathcal{V}$ for a complete or cocomplete
monoidal category $\mathcal{V}$, and the functor $T$ is a $\mathcal{V}$-functor.
The definitions transpose verbatim:

\begin{definition}[Enriched end and coend]
\label{def:enriched-end-coend}
For a $\mathcal{V}$-functor $T : \mathcal{C}^{\mathrm{op}} \otimes \mathcal{C}
\to \mathcal{V}$ (where $\mathcal{C}^{\mathrm{op}} \otimes \mathcal{C}$ is
the tensor product of $\mathcal{V}$-categories), the \term{enriched end}
$\int_c T(c, c)$ is the equalizer in $\mathcal{V}$ of two maps
$\prod_c T(c, c) \to \prod_{c, c'} \mathcal{V}(\mathcal{C}(c, c'), T(c, c'))$
encoding the wedge condition in $\mathcal{V}$. Dually for coend.
\end{definition}

\begin{proposition}[Enriched natural transformations as enriched end]
For $\mathcal{V}$-functors $G, H : \mathcal{C} \to \mathcal{V}$,
\begin{equation*}
  [\mathcal{C}, \mathcal{V}]_\mathcal{V}(G, H) \;=\; \int_c \mathcal{V}(G c, H c).
\end{equation*}
\end{proposition}

This is the statement that closed (and complete enough) Chapter~20 made
without proof; the foundation is now in place.

\subsection{Reorganizing}

In an enriched setting, the coend calculus rules — Fubini and Yoneda
reduction — continue to hold with the same proofs, replacing ``set
products'' by ``$\mathcal{V}$-tensors'' and ``set sums'' by
``$\mathcal{V}$-cotensors''. The ninja Yoneda lemma in enriched form reads
\begin{equation*}
  \int^c \mathcal{C}(c, c') \otimes F(c) \;\cong\; F(c')
\end{equation*}
in $\mathcal{V}$, where $\otimes$ is the $\mathcal{V}$-tensor.


% ═══════════════════════════════════════════════════════════════════════════
\section{Exercises}
\label{sec:end-coend-exercises}
% ═══════════════════════════════════════════════════════════════════════════

\begin{exercisesbox}
\begin{qa}
\qaitem{What is a dinatural transformation $\alpha : T \xRightarrow{..} S$?}{A family $\alpha_c : T(c,c) \to S(c,c)$ satisfying a hexagonal coherence condition involving both $T(f, \mathrm{id})$ and $T(\mathrm{id}, f)$ for each $f : c \to c'$.}
\qaitem{What is a wedge from $d$ to $T$?}{A dinatural transformation $\Delta d \xRightarrow{..} T$, equivalently a family $\omega_c : d \to T(c, c)$ satisfying $T(\mathrm{id}, f) \omega_c = T(f, \mathrm{id}) \omega_{c'}$ for every $f : c \to c'$.}
\qaitem{Define the end of $T : \mathcal{C}^{\mathrm{op}} \times \mathcal{C} \to \mathcal{D}$.}{An object $\int_c T(c,c) \in \mathcal{D}$ together with a universal wedge $\pi : \Delta \int_c T(c,c) \xRightarrow{..} T$.}
\qaitem{Express the end as an equalizer.}{$\int_c T(c,c) = \mathrm{eq}(\prod_c T(c,c) \rightrightarrows \prod_{f : c \to c'} T(c, c'))$, where the two arrows are induced by $T(\mathrm{id}, f)$ and $T(f, \mathrm{id})$.}
\qaitem{When does an end exist?}{When $\mathcal{D}$ is complete (has small limits) and $\mathcal{C}$ is small. Ends are limits.}
\qaitem{Show that natural transformations form an end.}{$\mathrm{Nat}(F, G) = \int_c \mathcal{D}(Fc, Gc)$: the wedge condition is the naturality square.}
\qaitem{What is the coend of $T$?}{An object $\int^c T(c,c)$ with universal cowedge $\iota : T \xRightarrow{..} \Delta \int^c T(c,c)$ — dual to the end.}
\qaitem{Express the coend as a coequalizer.}{$\int^c T(c,c) = \mathrm{coeq}(\coprod_{f : c \to c'} T(c', c) \rightrightarrows \coprod_c T(c, c))$.}
\qaitem{What is the tensor-product picture of a coend in $\mathbf{Set}$?}{$\int^c P(c) \times Q(c) \cong (\coprod_c P(c) \times Q(c))/\sim$, where $(P(f)(p), q) \sim (p, Q(f)(q))$ for $f : c \to c'$.}
\qaitem{Write the profunctor composition formula.}{$(Q \circ P)(c, a) = \int^b Q(c, b) \times P(b, a)$ for $P : \mathcal{B}^{\mathrm{op}} \times \mathcal{A} \to \mathbf{Set}$, $Q : \mathcal{C}^{\mathrm{op}} \times \mathcal{B} \to \mathbf{Set}$.}
\qaitem{Write the left Kan extension formula as a coend.}{$(\operatorname{Lan}_K F)(d) = \int^c \mathcal{D}(Kc, d) \cdot F(c)$, where $S \cdot X$ is the copower.}
\qaitem{State the Fubini theorem for coends.}{$\int^{(c,d)} T = \int^c \int^d T = \int^d \int^c T$, when the coends exist.}
\qaitem{State the Yoneda reduction (ninja Yoneda).}{$\int^c \mathcal{C}(c, c') \times F(c) \cong F(c')$ for $F : \mathcal{C} \to \mathbf{Set}$.}
\qaitem{Why is the ninja Yoneda lemma equivalent to ordinary Yoneda?}{Both express the same fact that $F$ is determined by its values at all $c$ via the representables $\mathcal{C}(c, -)$; one as an end (Hom-out of representables), the other as a coend (tensor with representables).}
\qaitem{State the co-Yoneda / density theorem in coend form.}{For $F : \mathcal{C} \to \mathbf{Set}$, $F(-) \cong \int^c \mathcal{C}(c, -) \cdot F(c)$. Every functor is a colimit of representables.}
\qaitem{What is the contravariant Yoneda reduction?}{$\int^c G(c) \times \mathcal{C}(c', c) \cong G(c')$ for $G : \mathcal{C}^{\mathrm{op}} \to \mathbf{Set}$.}
\qaitem{Why does the hexagonal condition collapse to a square for wedges?}{Because the constant functor $\Delta d$ has trivial action, so $\Delta d(f, \mathrm{id}) = \Delta d(\mathrm{id}, f) = \mathrm{id}_d$; the hexagon's two ``$\Delta d$ sides'' become the identity.}
\qaitem{What is the center of a category as an end?}{$Z(\mathcal{C}) = \int_c \mathcal{C}(c, c)$, the monoid of natural endomorphisms of $\mathrm{Id}_\mathcal{C}$.}
\qaitem{Geometric realization as a coend?}{$|X| = \int^{[n] \in \Delta} X_n \cdot \Delta^n$, where $X : \Delta^{\mathrm{op}} \to \mathbf{Set}$, $\Delta^n$ is the topological $n$-simplex, and $\cdot$ is the copower.}
\qaitem{Define the enriched end.}{For $T : \mathcal{C}^{\mathrm{op}} \otimes \mathcal{C} \to \mathcal{V}$, the equalizer in $\mathcal{V}$ of $\prod_c T(c,c) \rightrightarrows \prod_{c,c'} \mathcal{V}(\mathcal{C}(c,c'), T(c,c'))$.}
\qaitem{What is the enriched object of $\mathcal{V}$-natural transformations?}{$[\mathcal{C}, \mathcal{V}]_\mathcal{V}(G, H) = \int_c \mathcal{V}(Gc, Hc)$ in $\mathcal{V}$.}
\qaitem{When are ordinary natural transformations the special case of dinatural ones?}{When $T$ and $S$ factor through the projection to $\mathcal{C}$ (i.e., are covariant), the hexagonal condition collapses to the naturality square.}
\qaitem{What is a cowedge?}{A dinatural transformation $T \xRightarrow{..} \Delta d$: a family $\nu_c : T(c, c) \to d$ with $\nu_c \circ T(f, \mathrm{id}) = \nu_{c'} \circ T(\mathrm{id}, f)$.}
\qaitem{What is the relationship between ends/coends and ordinary limits/colimits?}{Ends are limits and coends are colimits of specific diagrams. The integral notation captures the ``contraction along the diagonal'' aspect that ordinary limit notation hides.}
\qaitem{Why does the Yoneda reduction fail without an integration variable?}{Because the reduction $\int^c \mathcal{C}(c, c') \cdot F(c) \cong F(c')$ is exactly the substitution of $c = c'$ that the integral notation indicates; without integration, the formula collapses to a single class and loses naturality.}
\qaitem{Coend version of currying for $\mathbf{Set}$-valued profunctors?}{$[\mathcal{A}^{\mathrm{op}}, [\mathcal{B}, \mathbf{Set}]] \simeq [\mathcal{A}^{\mathrm{op}} \times \mathcal{B}, \mathbf{Set}]$ — both expressible via coends/ends in symmetric ways.}
\qaitem{Why is Fubini for coends ``the same'' as Fubini for double colimits?}{Because a coend is the colimit of a specific double diagram, and Fubini for colimits says colimits over $\mathcal{C} \times \mathcal{D}$ may be computed iteratively.}
\qaitem{Write the coend formula for ``functor evaluation at $a$'' using representables.}{$F(a) \cong \int^c \mathcal{C}(c, a) \cdot F(c)$ — special case of Yoneda reduction with $c' = a$.}
\qaitem{If $\mathcal{C}$ is discrete, what is $\int^c F(c) \times G(c)$?}{$\coprod_c F(c) \times G(c)$ — no morphisms means no identifications, so the coend reduces to the coproduct.}
\qaitem{Why is the coend the ``right'' way to write profunctor composition?}{Because both projections from $T(c, c) = Q(c, b) \times P(b, a)$ — the contravariant action of $b$ on $Q$ and the covariant action of $b$ on $P$ — must be identified, exactly the coend's quotient relation.}
\end{qa}
\end{exercisesbox}


% ═══════════════════════════════════════════════════════════════════════════
\section*{Chapter Summary}
\addcontentsline{toc}{section}{Chapter Summary}
% ═══════════════════════════════════════════════════════════════════════════

\begin{summarybox}
\textbf{Dinatural transformations.}\quad For mixed-variance functors $T, S :
\mathcal{C}^{\mathrm{op}} \times \mathcal{C} \to \mathcal{D}$, a family
$\alpha_c : T(c,c) \to S(c,c)$ satisfying a hexagonal coherence. The constant
case collapses to wedges and cowedges — the dinatural analogues of cones.

\textbf{End.}\quad $\int_c T(c, c)$ is the universal wedge to $T$; equivalently,
the equalizer of two maps from $\prod_c T(c, c)$ to $\prod_{c \to c'} T(c, c')$.
Natural transformations are ends: $\mathrm{Nat}(F, G) = \int_c \mathcal{D}(Fc, Gc)$.

\textbf{Coend.}\quad $\int^c T(c, c)$ is the universal cowedge from $T$;
equivalently, the coequalizer of two maps to $\coprod_c T(c, c)$. Profunctor
composition and left Kan extensions are coends.

\textbf{The coend calculus.}\quad Fubini for ends and coends (swap order of
integration); Yoneda reduction (the ``ninja Yoneda lemma'':
$\int^c \mathcal{C}(c, c') \times F(c) \cong F(c')$); the co-Yoneda /
density theorem (every functor is a colimit of representables).

\textbf{Enriched ends and coends.}\quad The same definitions transpose to
$\mathcal{V}$-functors, with $\mathcal{V}$-products and $\mathcal{V}$-coproducts
replacing set products and sums. The enriched natural transformations object
of Chapter~20 is the enriched end $\int_c \mathcal{V}(Gc, Hc)$.
\end{summarybox}


% ═══════════════════════════════════════════════════════════════════════════
\section*{Formal Restatement}
\addcontentsline{toc}{section}{Formal Restatement}
% ═══════════════════════════════════════════════════════════════════════════

\begin{formalrestatement}
\textbf{Dinatural transformation $T \xRightarrow{..} S$.}\quad Family $\alpha_c
: T(c,c) \to S(c,c)$ with hexagonal coherence: $S(\mathrm{id}, f) \alpha_c
T(f, \mathrm{id}) = S(f, \mathrm{id}) \alpha_{c'} T(\mathrm{id}, f)$ for each
$f : c \to c'$.

\medskip
\textbf{End.}\quad
\begin{equation*}
  \int_c T(c, c) \;=\; \mathrm{eq}\bigl(
    \textstyle\prod_c T(c, c) \rightrightarrows
    \textstyle\prod_{f : c \to c'} T(c, c')
  \bigr).
\end{equation*}

\medskip
\textbf{Coend.}\quad
\begin{equation*}
  \int^c T(c, c) \;=\; \mathrm{coeq}\bigl(
    \textstyle\coprod_{f : c \to c'} T(c', c) \rightrightarrows
    \textstyle\coprod_c T(c, c)
  \bigr).
\end{equation*}

\medskip
\textbf{Natural transformations as end.}\quad
$\mathrm{Nat}(F, G) = \int_c \mathcal{D}(Fc, Gc)$.

\medskip
\textbf{Fubini.}\quad
$\int_{(c,d)} T = \int_c \int_d T = \int_d \int_c T$; coend version dual.

\medskip
\textbf{Yoneda reduction (ninja Yoneda).}\quad
$\int^c \mathcal{C}(c, c') \times F(c) \cong F(c')$ for $F : \mathcal{C} \to
\mathbf{Set}$; dually, $\int^c G(c) \times \mathcal{C}(c', c) \cong G(c')$
for $G : \mathcal{C}^{\mathrm{op}} \to \mathbf{Set}$.

\medskip
\textbf{Co-Yoneda / density.}\quad
$F \cong \int^c \mathcal{C}(c, -) \cdot F(c)$.

\medskip
\textbf{Profunctor composition.}\quad
$(Q \circ P)(c, a) = \int^b Q(c, b) \times P(b, a)$.

\medskip
\textbf{Left Kan extension formula.}\quad
$(\operatorname{Lan}_K F)(d) = \int^c \mathcal{D}(Kc, d) \cdot F(c)$.

\medskip
\textbf{Enriched end.}\quad For $T : \mathcal{C}^{\mathrm{op}} \otimes \mathcal{C}
\to \mathcal{V}$, $\int_c T(c, c) \in \mathcal{V}$ defined via equalizer of
$\mathcal{V}$-products encoding the $\mathcal{V}$-wedge condition.

\medskip
\noindent\textit{Chapter~27 uses the coend calculus to develop \emph{weighted
limits and colimits} in full generality. The Yoneda reduction is the
calculation tool that powers most of the enriched-categorical theory to
follow.}
\end{formalrestatement}
